\section{Introduction}
\label{sec:introduction}
Many applications require a choice among a finite set of answers, such as a
category, a rating, or a control action. A language model can make such a choice
by generating text, but it can also evaluate the candidate answers directly.
Direct evaluation avoids answer parsing and restricts the result to the allowed
set.

Dedicated decision models such as Jev use specialized training for these
tasks~\cite{jev-introduction}. We examine a complementary approach: efficient
decision inference from existing pretrained models, without requiring additional training.
JET (Justification Evaluation in Transformer) scores each answer by its
conditional likelihood and converts the scores into a choice, a binary
probability, or an ordinal estimate.

The contribution is an inference system for local execution. JET shares the
question prefix across candidates and isolates recurrent state during candidate
evaluation. These methods
reduce repeated work while retaining the likelihood-based decision rule.
CPU and GPU placement allow models to run under different memory constraints.
An optional task-aligned LoRA experiment tests whether a smaller model can
improve its decisions without changing the inference interface.

We evaluate decision accuracy and throughput jointly on MMLU and BoolQ using
consumer hardware. The comparison covers model choice and GPU hardware,
and execution optimizations, with Jev as a reference
(Figure~\ref{fig:mmlu-throughput}). Full-MMLU evaluation measures decision
quality and complete-process throughput; a controlled experiment on a fixed
subset measures execution improvements.
Section~\ref{sec:finetuning} reports the separate fine-tuning experiment and
its deployment limitations.
The observed 2.18--2.23-fold speedups preserve outputs in the tested comparison. These results assess the
quality and cost of local decisions; differing evaluation and timing protocols
do not establish parity with Jev.
